\section{Taxonomic Framework: The Five-Dimensional Constraint Lifecycle}
\label{sec:framework}

In renewable energy forecasting, designating a neural network as ``physics-constrained'' often groups disparate mathematical, structural, and empirical operations under a single umbrella term \cite{lipton2019}. To establish an operational foundation for systematic auditing, we define a structured \textbf{Constraint Lifecycle Taxonomy}.

\subsection{Unit of Analysis and Lifecycle Representation}
In our scoping framework, the primary unit of analysis is the individual \textbf{forecasting study} ($N=46$ unique scoping studies), characterized by its primary operational constraint mechanism. When a study incorporates multiple physical elements (e.g., combining an empirical power curve regularizer with a spatial wake graph architecture), it is classified according to the dominant physical constraint governing the computational graph, while auxiliary components are documented in the supplementary ledger. Formally, we characterize the primary constraint architecture $C$ of each audited study as a 5-tuple operating across orthogonal lifecycle stages:
\begin{equation}
C = \big\langle \mathcal{D}_1, \mathcal{D}_2, \mathcal{D}_3, \mathcal{D}_4, \mathcal{D}_5 \big\rangle,
\label{eq:constraint_tuple}
\end{equation}
where $\mathcal{D}_1$ represents parameter provenance, $\mathcal{D}_2$ denotes application stage, $\mathcal{D}_3$ specifies inference telemetry requirements, $\mathcal{D}_4$ captures implementation software fidelity, and $\mathcal{D}_5$ evaluates ablation semantics and harm reporting. The lifecycle trajectory from physical specification to operational evaluation is illustrated in Figure~\ref{fig:lifecycle}.

\begin{figure*}[!t]
\centering
\includegraphics[width=\textwidth]{fig_constraint_lifecycle.pdf}
\caption{The Five-Dimensional Constraint Lifecycle Framework. Tracking domain physical constraints from parameter origin ($\mathcal{D}_1$) through computational graph formulation ($\mathcal{D}_2$), operational telemetry dependency ($\mathcal{D}_3$), software implementation parity ($\mathcal{D}_4$), and empirical ablation validity ($\mathcal{D}_5$).}
\label{fig:lifecycle}
\end{figure*}

\subsection{Operational Definitions of Lifecycle Dimensions}
The taxonomy evaluates each audited constraint architecture across five orthogonal dimensions:

\paragraph{Dimension 1: Parameter Provenance ($\mathcal{D}_1$)}
Parameter provenance identifies the epistemic origin and calibration basis of the numerical values governing a constraint:
\begin{itemize}[leftmargin=1.5em,itemsep=1pt,topsep=2pt]
\item \textit{Intrinsic constraints}: parameters derive analytically from governing physical invariants, manufacturer nameplate limits, or classical equations without empirical fitting (e.g., the Betz limit $C_p \le 16/27$ or kinematic wake equations).
\item \textit{Historical-static priors}: parameters or reference distributions are estimated offline from historical operational SCADA archives and held fixed during inference (e.g., empirical power curves fitted via kernel density estimation).
\item \textit{Learned-online constraints}: parameters adapt dynamically during test-time inference via an auxiliary neural subnetwork.
\end{itemize}
When a study omits calibration details, it is designated as \textit{unreported}. Distinguishing provenance is essential: static empirical priors assume aerodynamic stationarity, risking silent miscalibration under turbine aging, blade degradation, or curtailment.

\paragraph{Dimension 2: Application Stage ($\mathcal{D}_2$)}
This dimension identifies where domain physics enters the computational pipeline:
(1)~\textit{training-loss regularization}, where physics acts as a differentiable penalty appended to the loss objective ($\mathcal{L} = \mathcal{L}_{\text{task}} + \lambda \mathcal{L}_{\text{phys}}$);
(2)~\textit{inference-architecture constraints}, which modify the network layers to project predictions onto admissible manifolds;
(3)~\textit{data-generation priors}, used upstream to synthesize training scenarios; and
(4)~\textit{post-hoc correction filters}, applied as external mathematical projections after inference.

\paragraph{Dimension 3: Inference Telemetry Requirements ($\mathcal{D}_3$)}
Inference telemetry identifies external data streams required during real-time deployment:
(1)~\textit{none}, relying solely on historical power measurements from the target turbine;
(2)~\textit{forecast-input}, requiring future meteorological inputs (e.g., forward numerical weather prediction wind vectors); and
(3)~\textit{contemporaneous-state}, requiring synchronized real-time measurements from adjacent turbines or met masts.
When a base forecasting model requires only past power series ($\mathcal{D}_3 = \texttt{none}$), offline training losses may nevertheless index empirical power curves against historical wind speed observations. In such setups, proxy staleness represents an offline training-loss sensitivity rather than an inference-time telemetry delay.

\paragraph{Dimension 4: Implementation Fidelity ($\mathcal{D}_4$)}
Implementation fidelity evaluates whether public software faithfully executes published mathematical formulations:
(1)~\textit{code-aligned}, where public source code strictly matches published equations;
(2)~\textit{code-mismatch}, where the codebase exhibits an algebraic or directional deviation from the text; and
(3)~\textit{not-found}, where no public repository or script accompanies the study.
Independent software inspection is necessary because published mathematical expressions can mask numerical discrepancies in computational graph execution.

\paragraph{Dimension 5: Harm Reporting and Ablation Semantics ($\mathcal{D}_5$)}
This dimension examines whether evaluations causally isolate physical constraints and report operational failure boundaries:
(1)~\textit{constraint-on-off}, where treatment and baseline share identical parameter capacity, differing only in the activation of the physical prior ($\lambda=0$ vs.\ $\lambda>0$);
(2)~\textit{module-ablation}, where the ablation excises an entire multi-parameter neural block, confounding structural capacity with domain knowledge;
(3)~\textit{both}, reporting both isolated toggles and architectural excisions; and
(4)~\textit{negative-effect-reported}, explicitly documenting operating regimes where physical constraints degrade downstream performance.

\subsection{Methodological Principles for Schema Disambiguation}
To ensure reproducible coding across independent raters, the framework enforces four disambiguation rules. Parameter provenance ($\mathcal{D}_1$) separates offline calibration from runtime telemetry ($\mathcal{D}_3$), ensuring that historical estimation windows are not conflated with active SCADA stream dependencies. Application stage ($\mathcal{D}_2$) is determined by the backpropagation path: loss regularizers update weights through objective penalties during training, whereas inference-architecture modifications constrain the forward pass under arbitrary parameter states. For composite methodologies, the primary classification is determined by the component directly modifying the forward computational graph, with secondary mechanisms logged in the ledger. Causal ablation validity ($\mathcal{D}_5$) requires parameter-matched controls; excising multi-layer neural submodules evaluates collective network capacity rather than isolating the specific inductive contribution of the physical constraint.
